\documentclass[11pt]{amsart}
\usepackage[margin=1in]{geometry}
\usepackage{amsmath,amssymb,mathtools,mathrsfs,microtype,enumitem,graphicx,needspace}
\usepackage{xr-hyper}
\usepackage[colorlinks=true,allcolors=blue]{hyperref}
\numberwithin{equation}{section}
\newtheorem{theorem}{Theorem}[section]
\newtheorem{proposition}[theorem]{Proposition}
\newtheorem{lemma}[theorem]{Lemma}
\newtheorem{corollary}[theorem]{Corollary}
\newtheorem{remark}[theorem]{Remark}
\newcommand{\Holder}[1]{C^{#1,1/2}}
\newcommand{\T}{\mathbb T^3}
\newcommand{\R}{\mathbb R}
\newcommand{\dd}{\,\mathrm d}
\newcommand{\Id}{\operatorname{Id}}
\newcommand{\curl}{\operatorname{curl}}
\newcommand{\diver}{\operatorname{div}}
\newcommand{\supp}{\operatorname{supp}}
\newcommand{\tr}{\operatorname{tr}}
\newcommand{\norm}[1]{\left\lVert#1\right\rVert}
\newcommand{\av}[1]{\langle #1\rangle_s}
\newcommand{\PP}{\mathbb P}
\newcommand{\ip}[2]{\langle #1,#2\rangle}
\newcommand{\eps}{\varepsilon}

\providecommand{\MainAuxPath}{../../../output/pdf/torus-rewrite/main}
\externaldocument[M-]{\MainAuxPath}[../torus-rewrite/main.pdf]
\title[Smooth terminal regularity]{Smooth terminal regularity away from the origin}
\author{}\date{}
\begin{document}
\raggedbottom
\begin{abstract}
The construction of Craik--Criminale waves on the torus can be chosen so that
the velocity extends smoothly in space and time through the singular time
away from the origin. We give the additional estimates for this strengthening
of the terminal conclusion in the main paper.
\end{abstract}
\maketitle

\section{Choice of \texorpdfstring{$\mathcal B$}{B} and the stronger conclusion}

We use the notation and the construction of
\emph{Finite-Time Singularity Formation for the 3-D Euler Equations}. References with prefix M refer to that paper.

\begin{theorem}\label{aux:terminal:smooth}
The function $\mathcal B$ of Subsection~M.\ref{M-cc:sec:parameters} may be chosen so that
the restriction of the solution in Theorem~M.\ref{M-cc:target:unforced} to
$[0,T_*)\times(\mathbb T^3\setminus\{0\})$ extends to a smooth field on an
open space-time neighborhood of
$\{T_*\}\times(\mathbb T^3\setminus\{0\})$.
All other conclusions of that theorem remain unchanged.
The Euler equations are imposed only for $0\le t<T_*$.
\end{theorem}

Subsection~M.\ref{M-cc:sec:parameters} admits any nondecreasing
function $\mathcal B$ that satisfies \textup{(M.\ref{M-cc:par:floor})}
and such that, for every integer $s\ge1$, $\mathcal B(s)$ bounds the
fixed-order constants of the paper at orders and depths at most
$7K(s)$. We choose one for which $\mathcal B(s)$ also bounds the
constants in the time-derivative estimates below with all derivative
orders and the depth at most $7K(s)$.
These quantities depend only on the orders and the depth:
the estimates below display dependence on the previous frequency parameter as fixed
powers and contain no constant depending on the constructed solution.
The function $\mathcal B$ is chosen before the initial frequency
parameter $N$, and the orders $q_j$ and the frequencies are defined by
\textup{(M.\ref{M-cc:uf:tower:eq:schedule})} with it.
Lemma~M.\ref{M-cc:uf:tower:lem:slow-orders} applies to this choice,
and the construction and the main theorem are obtained with it. Thus
the orders satisfy $q_j\le q_{j+1}\le q_j+1$ and $q_j\to\infty$, and
\begin{equation}\label{aux:terminal:growth}
 \mathcal B(q_{j+1})\le\lambda_{j-1}.
\end{equation}
The additional estimates concern the fields $U_j$ of this construction.

As in Section~M.\ref{M-cc:uf:sec:tower}, $O_j$ is the oscillation of
step $j$ and $R_j$ is the residual of $U_j$. The fields of step $j$ are
defined on $[0,S_j]$, and $S_j>T_\infty=T_*$.

\begin{lemma}\label{aux:terminal:mixed}
For all integers $m,q\ge0$ and all sufficiently large $j$,
\begin{equation}\label{aux:terminal:increment-mixed}
 \|\partial_t^q(U_j-U_{j-1}-O_j)\|_{C([0,S_j];\Holder{m})}
 \le\lambda_j^{-1}.
\end{equation}
The time derivatives are taken at fixed spatial point.
The starting index may depend on $(m,q)$.
\end{lemma}

By \textup{(M.\ref{M-cc:uf:tower:eq:leray})} the time derivatives of
$U_j$ involve those of $R_j$; the next lemma bounds both. The residual
of $U_j$ is built in step $j$ from the parent $U_{j-1}$, and at time
order $q$ it uses the residual of that parent only to time order $q-1$.
This is why the induction on the time order closes.

\begin{lemma}\label{aux:terminal:time}
For every integer $q\ge0$ there are an exponent $\psi_q$ and an index
$j_q$, depending only on $q$ and the fixed construction, such that for
$j\ge j_q$ and $0\le i\le q$,
\begin{equation}\label{aux:terminal:time-bounds}
 \|\partial_t^iU_j\|_{C([0,S_j];L^\infty)}\le\lambda_j^{\psi_q},\qquad
 \|\partial_t^iR_j\|_{C([0,S_j];L^\infty)}\le\lambda_{j+1}^{-5/2}.
\end{equation}
\end{lemma}

The fields $\partial_t^iU_j$ and $\partial_t^iR_j$ are supported on frequencies at most
$2\Lambda_j$ and $4\Lambda_j$. Thus
\textup{(M.\ref{M-cc:uf:trunc:eq:lp})} and
\eqref{aux:terminal:time-bounds} bound every spatial derivative of these
fields: a derivative of order $m$ costs
$C_m\Lambda_j^m=C_m\lambda_j^{5m}$.

\section{Time derivatives of the increments}
\begin{proof}[Proof of Lemmas~\ref{aux:terminal:mixed} and~\ref{aux:terminal:time}]
We argue by induction on $q$. For $q=0$,
Lemma~\ref{aux:terminal:mixed} is the second bound of
\textup{(M.\ref{M-cc:uf:tower:eq:increment})}, which holds as soon as
$q_j\ge m$,
and \eqref{aux:terminal:time-bounds} is \textup{(T2)} of Theorem~M.\ref{M-cc:uf:tower:thm:finite}, with $\psi_0=1$ and
$j_0=1$. Let $q\ge1$, and assume Lemma~\ref{aux:terminal:time} with
$q-1$ in place of $q$. We prove Lemma~\ref{aux:terminal:mixed} for this
$q$ and every $m$, and then Lemma~\ref{aux:terminal:time} for $q$.

Let $j-1\ge j_{q-1}$, and write $\mu=\lambda_{j-1}$ and $\lambda=\lambda_j$.
Put $T=S_j$, $U=U_{j-1}$, $R_U=R_{j-1}$,
$\Lambda_U=\Lambda_{j-1}=\mu^5$ and $h=\lambda^{-1/2}$.
Below, all exponents depend only on $q$, on $m$ once it
enters, and on the fixed construction; constants $C_q$ may depend also
on the label and phase order of the estimate in which they stand. An
\emph{old-parent polynomial} is a
bound $C_q\mu^{a_q}$; here $a_q$ denotes an exponent of this kind, which,
like $C_q$, may change from one estimate to the next. Each statement is for $j$ large
in terms of these, which raises only the starting index, since
$K_j$ tends to infinity. We use the step of
Subsection~M.\ref{M-cc:uf:sec:parent-step} as it is applied in
Section~M.\ref{M-cc:uf:sec:tower}: the parent $U$,
with flow $X$ and $M=(\nabla U)\circ X$; the depth $P$ of
\textup{(M.\ref{M-cc:uf:tower:eq:step-orders})}; and the norms of
Appendix~M.\ref{M-cc:uf:finite:sec:calculus}, at the label length
$\ell_*$ and the phase length $h$. There
$\mathsf D=\mu^9$ and
$T_\infty<T\le1$. Some constants below depend also on the depth $P$;
they stand in estimates used at orders at most $P+q+2$, which
is at most $7K(q_{j+1})$ for $j$ large. The
value $\mathcal B(q_{j+1})$ in
\eqref{aux:terminal:growth} bounds these and is at most $\mu$, so a
bound with such a constant is again an old-parent polynomial.
By the definitions of Section~M.\ref{M-cc:uf:sec:tower},
$T=S_j\ge2\tau_j\ge\mu^{-1/8}$. The numbers $T^{-1}$, $\ell_*^{-1}$
and every power of $\mathsf D$ with an absolute exponent are therefore
old-parent polynomials.

\emph{Time derivatives of the parent.}
By the assumed bounds \eqref{aux:terminal:time-bounds} for $U_{j-1}$,
\[
 \|\partial_t^iU\|_\infty\le\mu^{\psi_{q-1}},\qquad
 \|\partial_t^iR_U\|_\infty\le\lambda^{-5/2}
 \qquad(0\le i\le q-1).
\]
Differentiate \textup{(M.\ref{M-cc:uf:tower:eq:leray})} for $U$:
\[
 \partial_t^qU
 =-\PP\diver\sum_{i=0}^{q-1}\binom{q-1}i
                 \partial_t^iU\otimes\partial_t^{q-1-i}U
  +\partial_t^{q-1}R_U .
\]
The fields on the right are supported on frequencies at most $4\Lambda_U$ and involve time
derivatives of order at most $q-1$. The periodic Leray projection
is bounded on $\Holder{0}$, so
\textup{(M.\ref{M-cc:uf:trunc:eq:lp})} bounds
$\|\partial_t^qU\|_\infty$ by an old-parent polynomial. By
\textup{(M.\ref{M-cc:uf:trunc:eq:lp})} again, for every $n\ge0$,
\begin{equation}\label{aux:terminal:parent-time}
 \begin{aligned}
 \sum_{|\beta|=n}\|\partial^\beta\partial_t^iU\|_{L^\infty}
   &\le C_{n,q}\mu^{a_q}\Lambda_U^n\qquad(0\le i\le q),\\
 \sum_{|\beta|=n}\|\partial^\beta\partial_t^iR_U\|_{L^\infty}
   &\le C_n\Lambda_U^n\lambda^{-5/2}\qquad(0\le i\le q-1).
 \end{aligned}
\end{equation}
No time derivative of $R_U$ of order $q$ is used below.

\emph{Time derivatives of the coefficients.}
With $D_t=\partial_t+U\cdot\nabla$, the
$\ell$-fold material derivatives $D_t^\ell\nabla U$, $\ell\le q$, and
$D_t^\ell R_U$, $\ell\le q-1$, are
finite sums of products of at most $\ell+1$ of the fields in
\eqref{aux:terminal:parent-time} and their spatial derivatives of order
at most $\ell+1$, with exactly one factor from the second line in the
case of $R_U$. They are supported on frequencies at most $4(q+1)\Lambda_U$.
By \textup{(M.\ref{M-cc:uf:trunc:eq:lp})}, with $(q+1)\Lambda_U$ in
place of $\Lambda$, their derivatives of order $m$
are at most $C_m(q+1)^m\Lambda_U^m$ times their suprema. Since
$\Lambda_U=\mu^5\le\mathsf D$, this is the hypothesis of
\textup{(M.\ref{M-cc:uf:trunc:eq:parent-composition})} with
$C_m'=C_m(q+1)^m$.
Since $\partial_t^\ell M=(D_t^\ell\nabla U)\circ X$, the composition
estimates \textup{(M.\ref{M-cc:uf:trunc:eq:parent-composition})} and
\textup{(M.\ref{M-cc:uf:hier:eq:coefficient-bound})} of
Proposition~M.\ref{M-cc:uf:trunc:prop:parent} bound the norm
$\|\partial_t^\ell M\|_{[n]}$, $\ell\le q$, at every order $n$ by an
old-parent polynomial whose exponent depends on $q$, not on $n$. In the
same way, since
$\partial_t^\ell(R_U\circ X)=(D_t^\ell R_U)\circ X$,
\begin{equation}\label{aux:terminal:residual-jets}
 \|\partial_t^\ell(R_U\circ X)\|_{[n]}\le C_q\mu^{a_q}\lambda^{-5/2}
 \qquad(0\le\ell\le q-1).
\end{equation}

The remaining time jets follow from $\partial_t D_aX=MD_aX$, $((D_aX)^{-1})_t=-(D_aX)^{-1}M$,
$(\kappa/\lambda)_t=-M^T\kappa/\lambda$ and the differentiated inverse identities of
reciprocal coefficients, by a number of operations depending only on $q$, with
Proposition~M.\ref{M-cc:uf:trunc:prop:parent} and the product
estimate of Lemma~M.\ref{M-cc:uf:trunc:lem:fields}. Thus
$\|\partial_t^\ell\mathsf C\|_{[n]}\le C_q\mu^{a_q}$ for $0\le\ell\le q$ and
for $\mathsf C$ among $(D_aX)^{\pm1}$, $M$, $\kappa/\lambda$,
$\lambda^2/|\kappa|^2$, $K=(D_aX)^TD_aX$ and $\mathcal D=2(D_aX)^T\partial_t D_aX$. These bounds hold
at every order $n$, so no label derivative is
spent on a time
derivative, and their exponents depend neither on the order, the level, the depth nor the
new frequency parameter. They use the time derivatives of $U$ through
order $q$ only. The construction is jointly smooth by
Proposition~M.\ref{M-cc:uf:hier:prop:identity}, so no time regularity of an unknown
is assumed.

\emph{The two inverses with time derivatives.} At each order $n$, put
$\mathcal J_l(f)=\sum_{i\le l}T^i\|\partial_t^if\|/i!$
with the norms $\|\cdot\|_{[n],\gamma}$ or $\|\cdot\|_{[n]}$. Time differentiation holds the label
$a$ and the independent phase $s$ fixed; oscillations are divided by
$\gamma$ afterwards, and means are unweighted. The time and phase
suprema remain inside the sums over label multi-indices and phase orders. For $0\le\ell\le q$,
extracting $\partial_t^\ell f$ from $\mathcal J_q(f)$ costs
$\ell!T^{-\ell}$. By \textup{(M.\ref{M-cc:uf:hier:eq:graph})},
$(\partial_t^\ell f)^\natural
=(\partial_t+U\cdot\nabla_x)^\ell f^\natural$.
The factors $1/i!$ turn the time Leibniz rule into a convolution, so
Lemma~M.\ref{M-cc:uf:trunc:lem:fields} holds for $\mathcal J_q$ with the same
costs, and coefficient multipliers cost old-parent
polynomials.

A mean $B=D_aX v$ with source $f_{\rm m}$ satisfies, by
\textup{(M.\ref{M-cc:uf:mean:strong})} with $g=(D_aX)^Tf_{\rm m}$ and $\mathcal T=\PP K$
on divergence-free label fields, $\mathcal Tv_t=\PP(g-\mathcal Dv)$. Hence
\[
 \mathcal T\partial_t^{l+1}v=\PP\Big[\partial_t^l g
 -\sum_{i=0}^l\binom li(\partial_t^i\mathcal D)\partial_t^{l-i}v
 -\sum_{i=1}^l\binom li(\partial_t^iK)\partial_t^{l+1-i}v\Big],
\]
whose right side contains $v$ and the source only to time order $l$.
By the last assertion of Theorem~M.\ref{M-cc:uf:mean:inverse}, the
constant $A$ of \textup{(M.\ref{M-cc:uf:mean:base-constant})} bounds the
inverse of $\mathcal T$ on the divergence-free fields of $\Holder{0}$ at each
time. Then \textup{(M.\ref{M-cc:uf:mean:thresholds})} is condition
\textup{(M.\ref{M-cc:uf:trunc:eq:slow})} for the equation
$\mathcal Tw=\PP F$ with the coefficient $K$, and
Lemma~M.\ref{M-cc:uf:trunc:lem:lifting} in its H\"older form inverts
$\mathcal T$ in the norms $\|\cdot\|_{[n],1/2}$ at every order, with a
constant that is an old-parent polynomial.
Starting from the bound
\textup{(M.\ref{M-cc:uf:mean:holder-bound})} for $v$, induction over
$l\le q-1$ and \textup{(M.\ref{M-cc:uf:trunc:eq:holder-sup})} give, with
the norms $\|\cdot\|_{[n]}$ on the left and $\|\cdot\|_{[n+1]}$ on the
right, and independently of the depth,
\begin{equation}\label{aux:terminal:mean-gain}
 \mathcal J_q(B)\le C_q\mu^{a_q}\mathcal J_{q-1}(f_{\rm m}).
\end{equation}
Thus the mean velocity at time order $q$ uses its source only to time
order $q-1$.
No boundary condition is imposed on time derivatives; the undifferentiated
solution keeps its actual initial trace. A transverse solution satisfies
$V_t=-(\Id-2P_\kappa)MV+(\Id-P_\kappa)f$; the same induction from
\textup{(M.\ref{M-cc:uf:hier:eq:trans-bound})} bounds its time derivatives
in $\mathcal J_q$ by an old-parent polynomial times $\mathcal J_q(f)$, and
\textup{(M.\ref{M-cc:uf:hier:eq:pressure-from-trans})} bounds its
pressure. On $[t_0,T]$ the leading
oscillation $A_1=g\mathbf aF_h$ solves this equation with $f=0$, and the cutoff of
Subsection~M.\ref{M-cc:uf:sec:later-wave} satisfies $|\widetilde\chi^{(l)}|\le C_lr^l$,
with the rate $r$ of the comparison,
and $r\le\mu$ by item 3 of the proof of
Theorem~M.\ref{M-cc:uf:tower:thm:finite}. With
\textup{(M.\ref{M-cc:uf:tower:eq:seed-norm})} and
\textup{(M.\ref{M-cc:uf:hier:eq:size})}, the same induction bounds
$A_1$, $V_1=\widetilde\chi A_1$, their pressures and
$R_{\widetilde\chi}=\widetilde\chi'(A_1+C[A_1])$ of \textup{(M.\ref{M-cc:uf:hier:eq:switch})}
in the weighted $\mathcal J_q$ by an old-parent polynomial times
$\lambda^{-5/4}$.

\emph{Sources and the expansion.} By \textup{(M.\ref{M-cc:uf:hier:eq:fast-mean})},
$\mathbf d\cdot B_i=\mathbf d\cdot C_i=0$ and the vanishing mixed phase means,
the mean source is
\[
 \bar f_p=-\diver_{\mathbf d}\sum_{i+i'=p}
              (\av{Y_i\otimes Y_{i'}}+B_i\otimes B_{i'})\qquad(p\ge3),
\]
where
$(\diver_{\mathbf d}(\mathbf u\otimes \mathbf w))_\ell=\mathbf d_i(\mathbf u_i\mathbf w_\ell)$,
and $\bar f_2$ is the same expression minus $R_U\circ X$, by
\textup{(M.\ref{M-cc:uf:hier:eq:source})}.
The time derivatives of the divergence term fall only on the fields and
on the old coefficients in $\mathbf d$; no phase derivative or new wave vector
appears. The mean sources are measured in $\mathcal J_{q-1}$, which
\eqref{aux:terminal:mean-gain} allows. The later sources contain
the means $B_p$ themselves, and neither $\bar f_p$ nor a mean pressure.
The oscillatory sources are differentiated in the expression for
$(\partial_t+M)C[V]+\mathbf d\mathsf p$ in the proof of
Lemma~M.\ref{M-cc:uf:hier:lem:relay}, with
\textup{(M.\ref{M-cc:uf:hier:eq:potential-time})}, and in \textup{(M.\ref{M-cc:uf:hier:eq:supp})},
which contain only coefficient multipliers, one label derivative, a
normalized inverse wave vector, the existing factor $\kappa\partial_s$ and
the fields; weights are carried as in
Subsection~M.\ref{M-cc:uf:sec:hierarchy}.

By \eqref{aux:terminal:residual-jets} and $T\le1$, the norm
$\mathcal J_{q-1}(R_U\circ X)$ is at most an old-parent polynomial
times $\lambda^{-5/2}$, which is
\textup{(M.\ref{M-cc:uf:hier:eq:parent-residual})} for these norms with
that polynomial as a factor.
The proof of Theorem~M.\ref{M-cc:uf:hier:thm:depth} uses the parent and
the wave only through the bounds
\textup{(M.\ref{M-cc:uf:hier:eq:coefficient-bound})},
\textup{(M.\ref{M-cc:uf:mean:bound})},
\textup{(M.\ref{M-cc:uf:hier:eq:trans-bound})},
\textup{(M.\ref{M-cc:uf:hier:eq:size})},
\textup{(M.\ref{M-cc:uf:hier:eq:parent-residual})},
\textup{(M.\ref{M-cc:uf:hier:eq:completion-bound})} and
\textup{(M.\ref{M-cc:uf:hier:eq:relay-bound})}, and through
Lemma~M.\ref{M-cc:uf:trunc:lem:fields}. By the preceding paragraphs
these hold in the norms $\mathcal J_q$, with the mean sources and the
residual of the parent measured in $\mathcal J_{q-1}$, when each power
of $\mathsf D$ on their right sides is replaced by an old-parent
polynomial and the constants depend also on $q$. The residual of the
parent enters that proof only through the mean source at level two,
whose bound has a power of $\mathsf D$; the polynomial that replaces
this power is taken at least as large as the one in the bound for the
residual. The proof uses the hypothesis
$\mathsf D^C\lambda^{-1/12}\le1$ of the theorem only to absorb powers
of $\mathsf D$. With an old-parent polynomial in their place the
hypothesis reads $C_q\mu^{a_q}\lambda^{-1/12}\le1$, which holds for
$j$ large, since $\lambda=\mu^{K_j}$. Since time convolution does not change the factor
$\lambda^{-2/3}$ per level, that proof therefore
applies in $\mathcal J_q$, at every order, to the profiles $V_p$, $C_p$,
$B_p$, the
oscillatory pressures and the sources $f_p+f_p^S$, and gives
\textup{(M.\ref{M-cc:uf:hier:eq:Vbound})}--\textup{(M.\ref{M-cc:uf:hier:eq:Cbound})}
with each power $\mathsf D^C$ replaced by an old-parent polynomial
$C_q\mu^{a_q}$. The mean pressures are
not used.
At the depth $P$, at each order $n$ and with constants $C_{n,P}$ that
depend also on $q$, the same proof bounds $W_{\rm osc}$ by
$C_{n,P}\mu^{a_q}\lambda^{-5/4}$ in the weighted $\mathcal J_q$, the mean
sum $\bar W$ by
$C_{n,P}\mu^{a_q}\lambda^{-5/2}$ in the
unweighted one, and the remainder $R_P$ in the unweighted
$\mathcal J_q$ by
$C_{n,P}\lambda^{-5/4}\lambda^{-2P/3}$. The term
$R_{\widetilde\chi}$ is at most an old-parent polynomial times
$\lambda^{-5/4}$ in the weighted $\mathcal J_q$, and it vanishes, with
its time derivatives, outside $[t_0,t_0+1/r]$.

\emph{Eulerian time derivatives.} By
$\partial_t=(\partial_t+U\cdot\nabla_x)-U\cdot\nabla_x$, the Eulerian
derivative $\partial_t^if^\natural$ of a profile $f$, for $i\le q$, is a
finite sum of products of factors $\partial_t^l\partial^\beta U$, with
$l\le q-1$ and $|\beta|\le q-1$, and one factor
$\partial_x^{\beta'}(\partial_t^lf)^\natural$ with $l+|\beta'|\le q$.
Lemma~M.\ref{M-cc:uf:hier:lem:physical} converts the bounds in
$\mathcal J_q$ to Eulerian $\Holder{n}$ bounds of the last factor,
spending $n+|\beta'|+1$ label and phase derivatives. Each of these
derivatives costs $C\mathsf D^C\lambda^{3/2}\le\lambda^2$ for an oscillatory
profile, and $C\mathsf D^C$, an old-parent polynomial, for a
phase-independent one.
The extraction costs $q!T^{-q}$, an old-parent polynomial, and the
factors of $U$ cost
$\Lambda_U=\mu^5$ for each derivative, by
\eqref{aux:terminal:parent-time}.

For the phase-independent mean sum and fixed $(m,q)$ the whole cost is
an old-parent polynomial. This preserves $\lambda^{-5/2}$: with a
constant and an exponent that depend also on $m$, and for $j$ large,
\begin{equation}\label{aux:terminal:mean-mixed}
 \|\partial_t^q\bar W^\natural\|_{C([0,S_j];\Holder{m})}
 \le C_q\mu^{a_q}\lambda^{-5/2}\le\tfrac12\lambda^{-1},
\end{equation}
the last step because $\lambda=\mu^{K_j}$ and $K_j\to\infty$.

For the approximate solution $V=U+W^\natural$ of step $j$ and its
residual $R_V=(R_P+R_{\widetilde\chi})^\natural$, the same conversion
gives, for $0\le i\le q$ and $0\le n\le P$,
\begin{equation}\label{aux:terminal:V-mixed}
 \|\partial_t^iV\|_{C_t\Holder{n}}\le\lambda^{2(n+1)+\psi'_q},\qquad
 \|\partial_t^iR_V\|_{C_t\Holder{1}}
 \le\lambda^{\psi'_q}\lambda_{j+1}^{-4},
\end{equation}
with an exponent $\psi'_q\ge0$ independent of $j$ and $n$; the constants,
which depend on $n$, $q$ and the depth, are bounded by
$\mathcal B(q_{j+1})$. Indeed, the parent contributes at most $\lambda^{n+2}$ to the
first bound, by \eqref{aux:terminal:parent-time} and
$K_j\to\infty$. In the conversion, Leibniz' rule distributes
the $n$ derivatives over the factors of a product, and on a factor of $U$
each of them costs $\Lambda_U\le\lambda$. For
the residual, the remainder $R_P$ carries
$\lambda^{-2P/3}=\lambda_{j+1}^{-4}$, and on $[t_0,t_0+1/r]$
the conversion of Lemma~M.\ref{M-cc:uf:hier:lem:physical} carries the
weight $\gamma\le Ce^{-L_j}=C\lambda_{j+1}^{-4}$, by
\textup{(M.\ref{M-cc:uf:tower:eq:switch-small})} and
\textup{(M.\ref{M-cc:uf:tower:eq:schedule})}.

\emph{The truncation difference.}
Let $v_\Lambda=(S_{\Lambda_j}-\Id)V$ be the truncation difference of
Proposition~M.\ref{M-cc:uf:step:thm:step} at step $j$.
We have $\partial_t^qv_\Lambda=(S_{\Lambda_j}-\Id)\partial_t^qV$.
By the tail bound \textup{(M.\ref{M-cc:uf:trunc:eq:tail})} at the
order $P$, whose constant is at most $\mu$, and by
$\Lambda_j=\lambda^5$ and \eqref{aux:terminal:V-mixed}, for fixed $m$,
\[
 \|\partial_t^qv_\Lambda\|_{C_t\Holder{m}}
 \le\mu\Lambda_j^{m-P}\lambda^{2(P+1)+\psi'_q}
 =\mu\lambda^{-3P+5m+2+\psi'_q}.
\]
Since $\mu\le\lambda$ and $P\to\infty$, this is at most
$\frac12\lambda^{-1}$ for $j$ large. Since $V=U_{j-1}+W^\natural$,
$W=W_{\rm osc}+\bar W$ and $U_j=S_{\Lambda_j}V$, we have
$U_j-U_{j-1}-O_j=\bar W^\natural+v_\Lambda$, so this bound and
\eqref{aux:terminal:mean-mixed} prove
\eqref{aux:terminal:increment-mixed}.

\emph{The residual and the new velocity.} It remains to prove
Lemma~\ref{aux:terminal:time} for $q$. Since $U_j=S_{\Lambda_j}V$,
Lemma~M.\ref{M-cc:uf:trunc:lem:cutoff} and
\eqref{aux:terminal:V-mixed} at $n=0$ give
$\|\partial_t^iU_j\|_\infty\le\lambda^{\psi'_q+3}$ for $i\le q$, which is
the first bound of \eqref{aux:terminal:time-bounds} with
$\psi_q=\psi'_q+3$. By
\textup{(M.\ref{M-cc:uf:trunc:eq:residual})}, with
$v=v_\Lambda$ and $S=S_{\Lambda_j}$,
\[
 R_j=\PP\bigl[SR_V
 +\diver\bigl(v\otimes U_j+V\otimes v+(\Id-S)(V\otimes V)\bigr)\bigr].
\]
Differentiate $i\le q$ times in time; the derivatives distribute over
the products. The first term is bounded by
\eqref{aux:terminal:V-mixed}. Every other term contains the operator
$\Id-S$ applied to $\partial_t^lV$ or to
$\partial_t^l(V\otimes V)$ with $l\le q$. By
\eqref{aux:terminal:V-mixed}, the tail bound at the order $P$ and
\textup{(M.\ref{M-cc:holder:product-growth})}, these
terms are at most
\[
 C_{q,P}\Lambda_j^{2-P}\lambda^{4(P+1)+2\psi'_q}
 \le\lambda^{15-P+2\psi'_q}
 =\lambda^{15+2\psi'_q}\lambda_{j+1}^{-6}
\]
in $\Holder{1}$, because $\Lambda_j=\lambda^5$,
$C_{q,P}\le\lambda$ and $P=6K_{j+1}$.
Altogether $\|\partial_t^iR_j\|_\infty
\le\lambda^{15+2\psi'_q}\lambda_{j+1}^{-4}$ for $i\le q$, and
$\lambda^{15+2\psi'_q}\le\lambda_{j+1}^{3/2}$ for $j$ large, since
$K_{j+1}\to\infty$. This proves the second bound
of \eqref{aux:terminal:time-bounds} at the time order $q$, and the
induction is complete.
\end{proof}

\section{Extension away from the origin}

\begin{proof}[Proof of Theorem~\ref{aux:terminal:smooth}]
By Theorem~M.\ref{M-cc:uf:tower:thm:finite}, on $[0,S_j]$ the
oscillation $O_j$ is supported in the ball of radius
$\lambda_{j-1}^{-1}$ about the origin, and these radii tend to zero.
Fix a compact set $K\subset\mathbb T^3\setminus\{0\}$.
Only finitely many $O_j$ meet a neighborhood of $K$.
For every fixed pair $(m,q)$, \eqref{aux:terminal:increment-mixed} and
$\sum_j\lambda_j^{-1}<\infty$ give a summable bound for the
remaining increments in $C^q([0,T_\infty];C^m(K))$.
The finitely many earlier terms are smooth there because
$S_j>T_\infty$. The sum therefore converges with every mixed derivative
up to $T_\infty$ away from zero. Its limit equals $u$ before $T_\infty$
by Proposition~M.\ref{M-cc:uf:terminal:agreement}.

Seeley's extension theorem \cite[p.\ 625]{Seeley}, applied componentwise
in boundary charts and patched with a locally finite spatial partition
of unity, extends the field smoothly through $T_\infty$ away from zero.
The main paper identifies $T_*=T_\infty$, which proves the theorem.
\end{proof}

\begin{thebibliography}{9}
\bibitem{Seeley} R. T. Seeley, Extension of $C^\infty$ functions defined
in a half space, \emph{Proc. Amer. Math. Soc.} \textbf{15} (1964), 625--626.
\href{https://doi.org/10.1090/S0002-9939-1964-0165392-8}{doi:10.1090/S0002-9939-1964-0165392-8}.
\end{thebibliography}
\end{document}
